\documentclass[a4paper,12pt]{article}
\usepackage[utf8]{inputenc}
\usepackage[T1]{fontenc}
\usepackage[ngerman]{babel}
\usepackage{amsmath, amssymb}
\usepackage{geometry}
\usepackage{parskip}
\usepackage{titlesec}
\usepackage{enumitem}

\geometry{left=3cm, right=3cm, top=2.5cm, bottom=2.5cm}
\titleformat{\section}{\large\bfseries}{}{0em}{}[\titlerule]

\newcommand{\gentry}[3]{%
  \noindent\textbf{#1\quad #2}\nopagebreak\\[0.2em]
  #3\par\smallskip
}

\begin{document}

\begin{center}
  {\LARGE\bfseries Glossar zu Lektion 3}\\[0.5em]
  {\large Lineare Algebra (Modul 61112)}
\end{center}

\vspace{0.8em}

% ======================================================================
\section*{3.1 \quad Bilinearformen}

\gentry{3.1.1}{Bilinearform}{%
  Sei $K$ ein Körper, $U$ ein $K$-Vektorraum. Eine Abbildung $\beta: U \times U \to K$ heißt \textbf{Bilinearform}, wenn sie linear in beiden Argumenten ist:
  \begin{enumerate}[label=(\alph*), topsep=2pt, itemsep=1pt]
    \item $\beta(au + u', v) = a\beta(u,v) + \beta(u',v)$,
    \item $\beta(u, av + v') = a\beta(u,v) + \beta(u,v')$.
  \end{enumerate}
}

\gentry{3.1.9}{Satz ($\mathrm{Bil}_K(U)$ ist Vektorraum)}{%
  $\mathrm{Bil}_K(U)$, die Menge aller Bilinearformen auf $U$, ist ein $K$-Vektorraum.
}

\gentry{3.1.10}{Matrixdarstellung}{%
  Sei $B = (b_1,\ldots,b_n)$ eine geordnete Basis von $U$. Die \textbf{Matrixdarstellung} von $\beta$ bzgl.\ $B$ ist $M_B(\beta) := (\beta(b_i, b_j))_{i,j} \in M_{n,n}(K)$.
}

\gentry{3.1.14}{Satz (Bilinearform durch Matrix bestimmt)}{%
  Eine Bilinearform $\beta$ ist durch ihre Matrixdarstellung $M_B(\beta)$ eindeutig bestimmt.
  $\dim \mathrm{Bil}_K(U) = n^2$.
}

\gentry{3.1.17}{Transformationsformel}{%
  Sei $S = M_{BB'}(\mathrm{id}_U)$ die Basiswechselmatrix. Dann gilt:
  \[
    M_{B'}(\beta) = S^\top \cdot M_B(\beta) \cdot S.
  \]
}

\gentry{3.1.22}{Kongruenz}{%
  Zwei Matrizen $A, B \in M_{n,n}(K)$ heißen \textbf{kongruent} ($A \sim_K B$), wenn es eine invertierbare Matrix $S$ gibt mit $B = S^\top A S$.
  Zwei Bilinearformen $\beta_1, \beta_2$ heißen kongruent, wenn ihre Matrixdarstellungen kongruent sind.
}

\gentry{3.1.37}{Regulär / ausgeartet}{%
  Eine Bilinearform $\beta$ auf $U$ heißt \textbf{regulär}, wenn die zugehörige lineare Abbildung $\beta^{(1)}: U \to U^*$, $u \mapsto \beta(u,-)$, ein Isomorphismus ist. Andernfalls heißt $\beta$ \textbf{ausgeartet}.
  Äquivalent: $\det(M_B(\beta)) \neq 0$.
}

\gentry{3.1.43}{Zerlegungssatz}{%
  Sei $\beta$ eine Bilinearform auf $U$ und $U_0 \subseteq U$ ein \textbf{regulärer Unterraum} (d.\,h.\ $\beta|_{U_0}$ regulär). Dann gilt $U = U_0 \oplus U_0^\perp$.
}

\gentry{3.1.44/45}{Rang und Nullraum}{%
  Der \textbf{Rang} $\mathrm{Rg}(\beta) := \mathrm{Rg}(M_B(\beta))$.
  Der \textbf{Nullraum} $\mathrm{Rad}(\beta) := \ker(\beta^{(1)}) = \{u \in U \mid \beta(u,v) = 0 \; \forall\, v\}$.
}

\gentry{3.1.47}{Dimensionsformel}{%
  $\dim U = \mathrm{Rg}(\beta) + \dim \mathrm{Rad}(\beta)$.
}

% ======================================================================
\section*{3.2 \quad Symmetrische Bilinearformen}

\gentry{3.2.1}{Symmetrische Bilinearform}{%
  $\beta$ heißt \textbf{symmetrisch}, wenn $\beta(u,v) = \beta(v,u)$ für alle $u, v \in U$.
  Äquivalent: $M_B(\beta)$ ist symmetrisch für jede Basis $B$.
}

\gentry{3.2.10}{Polarisationsformel}{%
  Für symmetrisches $\beta$ (mit $\mathrm{char}(K) \neq 2$):
  $\beta(u,v) = \tfrac{1}{2}\bigl(\beta(u+v,u+v) - \beta(u,u) - \beta(v,v)\bigr)$.
}

\gentry{3.2.11}{Diagonalisierungssatz}{%
  Sei $\beta$ eine symmetrische Bilinearform auf dem endlich-dimensionalen $K$-Vektorraum $U$ mit $\mathrm{char}(K) \neq 2$. Dann gibt es eine Basis $B'$, sodass $M_{B'}(\beta)$ eine Diagonalmatrix ist.
}

\gentry{3.2.17/19}{Klassifikation über $\mathbb{C}$}{%
  Jede symmetrische Bilinearform über $\mathbb{C}$ mit Rang $r$ ist kongruent zur Form mit Matrix $\mathrm{diag}(1,\ldots,1,0,\ldots,0)$ ($r$ Einsen).
}

\gentry{3.2.25}{Trägheitssatz (Sylvester)}{%
  Jede reelle symmetrische Bilinearform mit Rang $r$ ist kongruent zu einer Diagonalform mit $p$ Einsen und $q = r-p$ negativen Einsen. Die \textbf{Signatur} $(p,q)$ ist unabhängig von der Wahl der Diagonalisierung.
}

\gentry{3.2.26}{Typ}{%
  Der \textbf{Typ} einer reellen symmetrischen Bilinearform ist die Signatur $(p,q)$. Zwei Formen sind kongruent genau dann, wenn sie denselben Typ haben.
}

\gentry{3.2.30}{Normalform über $\mathbb{R}$}{%
  Die Normalform ist $\mathrm{diag}(\underbrace{1,\ldots,1}_{p},\underbrace{-1,\ldots,-1}_{q},\underbrace{0,\ldots,0}_{n-r})$.
}

\gentry{}{Positiv/negativ (semi)definit}{%
  \begin{itemize}[topsep=2pt,itemsep=1pt]
    \item \textbf{positiv definit}: $\beta(u,u) > 0$ für alle $u \neq 0$ ($q = 0$, $r = n$).
    \item \textbf{negativ definit}: $\beta(u,u) < 0$ für alle $u \neq 0$ ($p = 0$, $r = n$).
    \item \textbf{positiv semidefinit}: $\beta(u,u) \geq 0$ für alle $u$ ($q = 0$).
    \item \textbf{indefinit}: $p > 0$ und $q > 0$.
  \end{itemize}
}

% ======================================================================
\section*{3.3 \quad Hermite'sche Formen}

\gentry{3.3.1}{Hermite'sche Form}{%
  Eine Abbildung $h: U \times U \to \mathbb{C}$ ($U$ ein $\mathbb{C}$-VR) heißt \textbf{Hermite'sche Form}, wenn:
  \begin{enumerate}[label=(\alph*), topsep=2pt, itemsep=1pt]
    \item $h(au + u', v) = \bar{a}\, h(u,v) + h(u',v)$ (konjugiert-linear in 1.\ Argument),
    \item $h(u, av + v') = a\, h(u,v) + h(u,v')$ (linear in 2.\ Argument),
    \item $h(v,u) = \overline{h(u,v)}$ (Hermite'sche Symmetrie).
  \end{enumerate}
  Insbesondere ist $h(u,u) \in \mathbb{R}$ für alle $u$.
}

\gentry{3.3.20}{Matrixdarstellung}{%
  Zu einer Basis $B$ ist $M_B(h) = (h(b_j, b_i))_{i,j}$. Diese Matrix ist \textbf{Hermitesch}: $M_B(h)^* = M_B(h)$, wobei $A^* := \overline{A^\top}$.
}

\gentry{3.3.25}{Transformationsformel}{%
  $M_{B'}(h) = S^* \cdot M_B(h) \cdot S$ mit $S^* = \overline{S^\top}$.
}

\gentry{3.3.43}{Trägheitssatz für Hermite'sche Formen}{%
  Jede Hermite'sche Form mit Rang $r$ hat eine eindeutige Signatur $(p,q)$ mit $p+q=r$. Die Normalform ist $\mathrm{diag}(1,\ldots,1,-1,\ldots,-1,0,\ldots,0)$.
}

\gentry{3.3.50}{Klassifikation}{%
  Zwei Hermite'sche Formen auf einem endlich-dimensionalen $\mathbb{C}$-Vektorraum sind genau dann kongruent, wenn sie dieselbe Signatur haben.
}

\end{document}
